\section{Interpretation and Use Limitations}
\label{sec:case-study-interpretation}

The representative case shows that Stockrium Lab can turn a supported configuration into an inspectable historical record. A user can see which inputs were selected, follow long-running backtests through persisted status, review symbol-level and aggregate outputs, open detailed visualizations, and compare compatible saved runs. Negative, mixed, or incomplete outcomes remain visible instead of being removed because they do not support a favorable narrative.

The appropriate use of these outputs is exploratory. A user may use them to identify differences between configurations, investigate unusual securities, review the evidence behind decisions, and decide what requires further verification. Return, win-rate, confidence, or risk fields are descriptive values produced under the saved simulator assumptions; none is a promise about a future transaction or a substitute for financial judgment.

\subsection{Limitations Visible to the User}
\label{subsec:case-study-limitations}

Table~\ref{tab:case-study-limitations} summarizes the most important cautions for interpreting the saved case. The list is intentionally concise because the objective is to explain the example, while Chapter~4 already documents the detailed simulation rules.

\begin{table}[H]
    \centering
    \caption{Principal limitations of the representative Stockrium Lab case}
    \label{tab:case-study-limitations}
    \footnotesize
    \renewcommand{\arraystretch}{1.2}
    \begin{tabularx}{\textwidth}{>{\raggedright\arraybackslash}p{3.5cm} X}
        \toprule
        \textbf{Limitation} & \textbf{Effect on interpretation} \\
        \midrule
        Current VN30 membership & The selected symbols are current constituents, not a survivorship-bias-free reconstruction of historical membership. \\
        \addlinespace
        Fundamental publication timing & Annual rows are bounded by financial year, but the database does not record the exact date on which every report became publicly available. \\
        \addlinespace
        Different comparison rules & The full pipeline and engine-only path differ in both decision acceptance and trade management, so their difference cannot be attributed to one component alone. \\
        \addlinespace
        Model and data reproducibility & Live sources and model-backed analysis may change; the configuration fingerprint does not guarantee identical future output. \\
        \addlinespace
        Simplified execution & Daily bars, fixed transaction costs, and per-symbol simulation omit spread, liquidity, price-limit queues, market impact, shared capital, and portfolio exposure. \\
        \addlinespace
        Background execution & FastAPI background tasks provide status polling but not durable resume, cancellation, or automatic retry after a process failure. \\
        \bottomrule
    \end{tabularx}
\end{table}

\subsection{Case Study Conclusion}
\label{subsec:case-study-conclusion}

This case is evidence of how a user can operate Stockrium Lab: configure a run, wait for completion, retain the result, reopen its details, and compare it with other outputs. The saved VN30 values demonstrate that the interface preserves trade-offs and unfavorable comparisons, including the stronger buy-and-hold result for most selected securities. The conclusion is therefore about tool capability and transparency, not investment effectiveness. Stronger market-performance claims would require separate point-in-time datasets, fixed model versions, preregistered comparisons, realistic portfolio accounting, and broader out-of-sample evidence.
